%============================================================
\chapter{Bai Hoc: Biet On Va Tri An (Gratitude and Appreciation)}
\begin{center}
  \textit{\color{truyenblue}Biet on mo ra su phong phu cua cuoc song; no bien nhung gi ta co thanh du day.}\\
  \textit{(Gratitude opens the fullness of life; it turns what we have into enough.)}\\[4pt]
  \small\color{truyengold}--- Thien nhien day ta ---
\end{center}
\ngancach

\section{Cho Trung Thanh Hachiko}
\begin{truyen}{Cho Trung Thanh Hachiko}{Cho Hachiko}
\chuhoa{H}{}H achiko moi ngay den ga tau Shibuya don chu ve nha sau gio lam viec tan ca.\\
\textit{(Hachiko came to Shibuya station every day to welcome his owner home after work ended.)}

Mot ngay chu qua doi dot ngot tai noi lam viec va khong bao gio tro ve nua.\\
\textit{(One day the owner died suddenly at work and never returned.)}

Nhung Hachiko van tro lai ga tau moi ngay mot cach deu dan trong muoi nam lien tiep.\\
\textit{(But Hachiko still returned to the station every day consistently for ten consecutive years.)}

Ong khong biet khai niem cai chet; ong chi biet chu minh la nguoi xung dang duoc cho.\\
\textit{(He did not understand the concept of death; he only knew his owner was someone worth waiting for.)}

Long biet on va trung thanh cua Hachiko khong the goi ten bang bat ky loi nao co du.\\
\textit{(Hachiko's gratitude and loyalty cannot be named by any words sufficient enough.)}

\end{truyen}
\begin{baihoc}
  Trung thanh la biet on theo thoi gian: ta co the biet on mot nguoi ca khi ho khong con o ben ta nua.\\
  \textit{(Loyalty is gratitude over time: we can be grateful to someone even when they are no longer beside us.)}
\end{baihoc}

\section{Ca Heo Biet On Nguoi Cuu}
\begin{truyen}{Ca Heo Biet On Nguoi Cuu}{Ca Heo}
\chuhoa{M}{}M ot nguoi tho lan thay mot con ca heo bi day luoi nguy hiem, linh so cat day giai cuu no.\\
\textit{(A diver saw a dolphin dangerously entangled in nets, instinctively cutting the ropes to free it.)}

Con ca heo boi ra tu do, voi lon vong quanh nguoi tho dan de vet thuong.\\
\textit{(The dolphin swam free, then made large circles around the diver checking for wounds.)}

No con sung vau vao tay nguoi de to ra su tiep xuc than thiet truoc khi boi di.\\
\textit{(It also nudged the person's hand to show intimate contact before swimming away.)}

Ba ngay sau con nguoi do gap kho khan khi lan sau; mot dan ca heo xuat hien day anh len bo.\\
\textit{(Three days later the same person got into difficulty diving deep; a dolphin pod appeared and pushed him to shore.)}

Lieu co the la con ca heo do quay lai hay khong, ta khong biet chac, nhung huong ve phia biet on thi co.\\
\textit{(Whether it was that same dolphin we cannot know for sure, but the direction of gratitude points there.)}

\end{truyen}
\begin{baihoc}
  Nhung gi ta trao di voi long thanh that co the quay tro lai theo nhieu cach ta khong lan tuong duoc.\\
  \textit{(What we give with true sincerity can return in ways we could never have imagined.)}
\end{baihoc}

\section{Suc Seo Che Cho Loai Khac}
\begin{truyen}{Suc Seo Che Cho Loai Khac}{Suc Seo}
\chuhoa{N}{}N guoi quan sat thay nhung con chim nho thuong xay to tren canh cay keo chan hom.\\
\textit{(Observers noticed small birds often build nests in the branches of acacia trees housing ants.)}

Luc dau nghi chim la nguyen nhan, nhung ky ra la chim chu dong tim den de suc nhung lo an tu cay.\\
\textit{(At first thought birds were the cause, but it turned out birds actively sought out trees with ant colonies.)}

Kien bao ve cay khoi sau benh rot vao cay va do lat loai chim an trung chim xuong bai.\\
\textit{(Ants protected the trees from pest borers and drove away thieving birds that stole nestling eggs.)}

Chim tu nguyen xay to o do de duoc bao ve mien phi tu he thong an ninh kien dau tien.\\
\textit{(Birds voluntarily nested there to receive free protection from the ant security system.)}

Su biet on duoc the hien bang viec tro thanh mot phan cua he sinh thai ho tro lan nhau.\\
\textit{(Gratitude is expressed by becoming part of a mutually supportive ecosystem.)}

\end{truyen}
\begin{baihoc}
  Khi ai do giup ban, ti le biet on tot nhat khong chi la loi cam on ma la tro nen co ich cho ho.\\
  \textit{(When someone helps you, the best expression of gratitude is not just saying thank you but becoming useful to them.)}
\end{baihoc}

\section{Con Rua Biet On Nguoi Giai Cuu}
\begin{truyen}{Con Rua Biet On Nguoi Giai Cuu}{Con Rua}
\chuhoa{T}{}T ruyen co ke ve con rua ca bien biet on nguoi fisherman da giai thoat no khoi luoi.\\
\textit{(The old tale tells of a sea turtle grateful to the fisherman who freed it from the net.)}

Rua quay lai dan nguoi ra bien sau con bao chi huong den ben bo an toan.\\
\textit{(The turtle returned and guided the person out to sea after the storm toward a safe harbor.)}

Phong tuc nguoi Nhat Boc Bong goi la urashima: biet on phai dam nhan tranh nhiem tra lai.\\
\textit{(The Japanese custom called urashima: gratitude must dare to take responsibility in return.)}

Cau chuyen co tich nen khong the kiem chung, nhung cam xuc ma no truyen dat la that.\\
\textit{(The fairy tale cannot be verified, but the feeling it conveys is real.)}

Biet on khong chi la cam giac; no la quyet dinh hanh dong tra lai cho cuoc song nhung gi cuoc song da trao.\\
\textit{(Gratitude is not just a feeling; it is the decision to act and give back to life what life has given.)}

\end{truyen}
\begin{baihoc}
  Biet on hoat dong nhu mot vong luon: cang trao di cang nhan lai, doi song cang tro nen phong phu.\\
  \textit{(Gratitude operates as a cycle: the more you give back the more you receive, life becomes ever richer.)}
\end{baihoc}

\section{Hoa Huong Duong Quay Mat Theo Nang}
\begin{truyen}{Hoa Huong Duong Quay Mat Theo Nang}{Hoa Huong Duong}
\chuhoa{M}{}M oi sang hoa huong duong huong mat ve phia dong noi mat troi moc sap den.\\
\textit{(Every morning sunflowers face east where the sun will rise.)}

Ca ngay chung xoay than theo hanh trinh cua mat troi tu dong sang tay tu tung phut.\\
\textit{(All day they slowly turn their stems following the sun's journey from east to west minute by minute.)}

Ban dem hoa quay tro lai vi tri dong de san sang chao mat troi sang hom sau.\\
\textit{(At night the flowers turn back to the east position to be ready to welcome the next day's sun.)}

Hanh dong xoay cua than hoa khong phai phan xa mac dinh; do la su the hien long biet on.\\
\textit{(The stem's turning movement is not just a default reflex; it is an expression of gratitude.)}

Biet on la huong mat ve noi soi sang cuoc song ban; lam dieu do moi ngay.\\
\textit{(Gratitude is facing toward what gives your life light; do that every day.)}

\end{truyen}
\begin{baihoc}
  Hay bat dau moi ngay bang viec quay mat ve phia nhung dieu va nhung nguoi dem anh sang vao doi ban.\\
  \textit{(Begin every day by turning your face toward the things and people who bring light into your life.)}
\end{baihoc}

\section{Cay Hoa Anh Dao Va Mua Xuan}
\begin{truyen}{Cay Hoa Anh Dao Va Mua Xuan}{Hoa Anh Dao}
\chuhoa{M}{}M oi nam cay hoa anh dao chi no trong mot tuan ngan ngui roi rung het trong nac long.\\
\textit{(Each year cherry blossom trees bloom for only one short week then fall completely in heartbreak.)}

Nguoi Nhat coi tuan no do la tang qua thieng lieng nhat cua thien nhien va tri an tung canh.\\
\textit{(The Japanese consider that blooming week the most sacred gift of nature and cherish every petal.)}

Ho ngoi ben duoi tan canh, khong dung dien thoai, chi lang nghe gio va nhin canh hoa bay.\\
\textit{(They sit beneath the canopy, not using phones, only listening to wind and watching petals fly.)}

Hanh dong do goi la mono no aware: cam nhan ve ve dep nang neu cua su qua do.\\
\textit{(That act is called mono no aware: feeling for the fleeting beauty of impermanence.)}

Biet on ve dep cua nhung dieu qua do la loai hinh tri tue tinh te nhat con nguoi co the phat trien.\\
\textit{(Appreciating the beauty of fleeting things is the most refined form of intelligence humans can develop.)}

\end{truyen}
\begin{baihoc}
  Biet on nhung gi dang co mat hom nay vi ngay mai no co the mat di khong thong bao truoc.\\
  \textit{(Be grateful for what is present today because tomorrow it may disappear without notice.)}
\end{baihoc}

\section{Vooc Cham Soc Nguoi Gia}
\begin{truyen}{Vooc Cham Soc Nguoi Gia}{Vooc}
\chuhoa{T}{}T rong bays vooc chau Phi, cac con tre cham soc va nhg tiep thuc an cho cac con gia yeu duoi.\\
\textit{(In African vervet monkey troops, young members care for and bring food to elderly weak members.)}

Vooc gia khong con co the leo hay hanh dong nhanh nhay nhung van duoc bao ve trong bays.\\
\textit{(Old vervets can no longer climb or act swiftly but are still protected within the troop.)}

Cac nha nghien cuu phat hien vooc gia cung cap tri thuc kinh nghiem quan trong cho ca bay.\\
\textit{(Researchers found elderly vervets provide important experiential knowledge for the whole troop.)}

Ho nho nhung nguon nuoc cu, nhung con duong an toan va nhung mua ba tung xay ra.\\
\textit{(They remember old water sources, safe paths, and past predator seasonal patterns.)}

Cham soc nguoi gia khong phai tu thien; do la dau tu khon ngoan vao kho luyen tri tue cua ca tap the.\\
\textit{(Caring for the elderly is not charity; it is a wise investment in the entire group's accumulated wisdom.)}

\end{truyen}
\begin{baihoc}
  Hay de nguoi lon tuoi trong nhom cua ban truyen dat tri thuc; ho la thu vien song cua tap the ban.\\
  \textit{(Let older people in your group share their knowledge; they are the living library of your community.)}
\end{baihoc}

\section{Ca Heo Giai Cuu Nguoi Bi Chuot Ran}
\begin{truyen}{Ca Heo Giai Cuu Nguoi Bi Chuot Ran}{Ca Heo}
\chuhoa{T}{}T ai New Zealand mot nguoi boi bi dan ca map vay quanh tren bien la ngoai bo.\\
\textit{(In New Zealand a swimmer was surrounded by sharks far offshore.)}

Mua dan ca heo bay den tao thanh vong tron bao ve xung quanh nguoi boi do.\\
\textit{(A pod of dolphins immediately arrived forming a protective circle around the swimmer.)}

Chung khong chuyen di ma bay lai trong 40 phut cho den khi tau cuu ho den.\\
\textit{(They did not leave but swam in circles for 40 minutes until the rescue boat arrived.)}

Mot ca heo con lien tuc vu vao nguoi boi de giup nguoi boi tinh tao khong ngo chim.\\
\textit{(One dolphin continuously nudged the swimmer to help them stay alert and not sink drowsy.)}

Loai ca heo khong co ly do gi de cuu con nguoi ngoai bat dong loai; nhung chung van cuu.\\
\textit{(Dolphins have no reason to rescue humans outside of interspecies empathy; yet they do.)}

\end{truyen}
\begin{baihoc}
  Long biet on va su giup do co the buoc qua ranh gioi loai; hay mo rong vong tay ban ra xa hon.\\
  \textit{(Gratitude and helping can cross the boundary of species; extend your circle of care wider.)}
\end{baihoc}

\section{Khi Con Cuu Minh Khoi Cuon}
\begin{truyen}{Khi Con Cuu Minh Khoi Cuon}{Con Khi}
\chuhoa{K}{}K hi lu cuon no lai, mot nha khoa hoc bi mat cu tren bo song ngan duoc cay lon.\\
\textit{(When flood waters swept and rose, a scientist was cut off on a riverbank blocked by a large tree.)}

Mot con khi nho bay den, trao cho anh mot doan roi day leo de ket vao canh cay moc danh kia.\\
\textit{(A small monkey came over, handing him a length of twisted vine to tie to the leaning tree branch.)}

Anh ket vao and dung dan bo day boi qua doln song lai an toan.\\
\textit{(He tied on and climbed the vine across the flooded river safely.)}

Khi day khoa hoc phan tich, loai khi do khong co loi ich cu the nao tu hanh dong nay het.\\
\textit{(When later scientists analyzed it, that monkey species had no concrete benefit from this action at all.)}

Altruism, su giup do vo tu, co the xuat hien o nhieu loai hon la chung ta dang nghi.\\
\textit{(Altruism, selfless helping, may appear in more species than we currently think.)}

\end{truyen}
\begin{baihoc}
  Du ban co la nguoi hay la thu vat, thien tuyen co the thuc day hanh dong tot dep khong tinh toan.\\
  \textit{(Whether you are human or animal, innate goodness can drive selfless acts without calculation.)}
\end{baihoc}

\section{Rua De Trung Tren Bai Bien}
\begin{truyen}{Rua De Trung Tren Bai Bien}{Rua Bien}
\chuhoa{T}{}T rong hanh trinh tuong tu mo cua minh, rua bien chon dung bai bien noi no chao doi de de trung.\\
\textit{(In a journey resembling a circle of memory, the sea turtle chooses the exact beach where it was born to lay eggs.)}

Sau hai muoi nam boi khap dai duong no quay lai noi bat dau chinh xac den tung met.\\
\textit{(After twenty years crossing the entire ocean it returns to its starting point accurate to the meter.)}

No khong biet gi ve bai bien khac; no chi biet noi day la nha va no phai tra lai.\\
\textit{(It knows nothing of other beaches; it only knows here is home and it must return.)}

Bai hoc la: biet on noi ban duoc sinh ra, noi da tao nen con nguoi ban hom nay.\\
\textit{(The lesson is: be grateful for where you were born, the place that created the person you are today.)}

Quay ve nguon coi de trao lai qua tang la hanh dong biet on dep de nhat ma sinh vat biet lam.\\
\textit{(Returning to your origins to give a gift back is the loveliest act of gratitude any creature knows how to do.)}

\end{truyen}
\begin{baihoc}
  Hay nho den noi ban bat dau, nhung nguoi dau tien giup ban, va tim cach tra lai dieu gi do co y nghia.\\
  \textit{(Remember where you started, the first people who helped you, and find a way to give back something meaningful.)}
\end{baihoc}
